\sectionguard
\section*{The Propers: Interpretive Possibilities}

\noindent\emph{The following six proposals are exploratory editorial
synthesis, not recovered compiler intent or teaching attributed to a cited
witness. Each names its closest checked precedent and the condition under which
the proposal remains sound.}

\Needspace{8\baselineskip}
\noindent\textbf{The ear governs the tongue.\enspace\properrefs{Ep., All., Gosp.}}
\emph{Classification: near analogue located.} Paul hands on what he first
received; the Corinthians receive, stand, hold fast, and are saved through that
Gospel. Psalm~80's larger context turns festival music into covenant hearing.
Christ then opens ears before releasing the tongue and makes the resulting
speech \latin{recte}; his command of silence prevents mere publicity from being
identified with confession. Together the elements propose a theology of
Christian speech: apostolic content precedes expression, the Spirit gives
interior hearing, and deed verifies utterance. Anthony §§13--14, Augustine's
\emph{De consensu} IV.4.5, Gregory I.10.20, and Honorius IV.68 supply partial
precedents. The proposal concerns the order of faith and witness; bodily
deafness is a condition Christ compassionately heals, not culpable unbelief.

\Needspace{8\baselineskip}
\noindent\textbf{Feared silence is answered by a groan and one word.\enspace\properrefs{Grad., Gosp., Off.}}
\emph{Classification: editorial liturgical synthesis; no direct precedent is
claimed.} The Gradual's
centonized versicle asks both \latin{ne sileas} and \latin{ne discedas a me}.
Christ answers not with extended discourse but by remaining with the sufferer,
touching him, looking heavenward, groaning, and speaking one effective
\latin{Ephphetha}. The Offertory then takes up the Gradual's
\latin{clamavi} and completes it with \latin{sanasti me}. The sequence suggests
that divine response is known by efficacious presence before it is measured by
verbal volume. Direct Psalm~27 and Mark~7 witnesses treat the components; the
proposal offers a spirituality of
persevering prayer, not a promise of an audible reply or a claim that every
experienced silence is abandonment.

\Needspace{8\baselineskip}
\noindent\textbf{Grace creates a gift capable of return.\enspace\properrefs{Coll., Ep., Sec., Comm.}}
\emph{Classification: near analogue located in an older formulary.} The
Collect's abundance exceeds merit and request; grace makes the persecutor an
apostle and does not remain void; the Secret names real service and gift while
asking God for acceptance and medicinal fruit; the Communion returns substance
and firstfruits to the Lord. The mechanism is neither passive reception nor
independent achievement. Grace gives identity and capacity; the creature acts;
the act returns Godward; divine acceptance makes it fruitful for the weak.
Honorius IV.65 joins Collect, Epistle, and Communion within an older Sunday set,
but not the present Secret. Augustine, Aquinas, Bede, and Salvian provide
doctrinal and moral analogues. The proposal remains bounded by the appointed
Pauline ending and may not smuggle in the unproclaimed remainder of v.~10.

\Needspace{8\baselineskip}
\noindent\textbf{The enemy may be conquered as a witness.\enspace\properrefs{Int., Coll., Ep.}}
\emph{Classification: partial saintly analogue located.} The Introit asks God
to arise and scatter his enemies. The Collect invokes mercy beyond merit. Paul
then names himself a persecutor to whom the risen Christ appeared and whom
grace made an apostle. Saint Oscar Romero explicitly joins Paul's transformation
to this Gospel's healing power: Christ remakes the sinner rather than
destroying him. As a relation within the Mass, the texts open a conversion-shaped
form of divine victory: enmity is scattered when the adversary is incorporated
into truthful testimony and service. No checked Psalm~67 commentator directly
joins the verse to Paul. Conversion does not erase judgment, repair owed, or
freedom; it reveals the paschal power by which an enemy can become a brother.

\Needspace{8\baselineskip}
\noindent\textbf{Material signs minister to the whole person.\enspace\properrefs{Gosp., Sec., Postcomm.}}
\emph{Classification: near analogue located.} Fingers, saliva, ears, tongue,
gaze, groan, and spoken command belong to Christ's healing act. At the altar a
material offering becomes accepted service and help for fragility; after
Communion the sacrament is asked to aid mind and body. The sequence proposes a
sacramental anthropology: created realities can serve personal grace because
the incarnate Lord acts through them, and salvation perfects rather than
abandons embodied nature. Ambrose's mystagogy, the Roman \emph{Ephpheta} rite,
Anthony, Irenaeus, Romero, and Schuster provide close analogues in distinct
registers. The historical miracle, preparatory baptismal sign, Eucharistic
sacrifice, and sacramental fruit remain distinct; their unity is Christ's care
for the one embodied person and the resurrection promised to that person.

\Needspace{8\baselineskip}
\noindent\textbf{From unvoid grace to heavenly fullness.\enspace\properrefs{Coll., Ep., Comm., Postcomm.}}
\emph{Classification: near analogue located.} The formulary develops a verbal
scale of abundance: divine goodness surpasses what prayer can measure; grace in
Paul is not \latin{vacua}; barns fill and presses overflow; the received
sacrament tends toward \latin{plenitudo caelestis remedii}. The intermediate
propers show how fullness travels: Gospel is received, senses are opened,
praise is returned, service is accepted, and firstfruits are given. Abundance
is therefore not stockpiled surplus but grace reaching apostolic, liturgical,
economic, and eschatological fruit. Honorius, Bede, Aquinas, and Schuster
provide partial analogues, not the entire chain. The proposal fails when
converted into prosperity teaching; its final measure is the common salvation
of mind and body, not private accumulation.
